Table~\ref{tab:main} and Table~\ref{tab:rank} give the main results; Fig.~\ref{fig:params} (which omits the from-scratch system) plots per-seed accuracy against trainable parameters. The pretrained model solves its own task (pretask accuracy 0.999) and none of the adapted tasks without adaptation (exact match 0.000 on descending sort, 0.001 on reversal).

\begin{table}[t]\centering\caption{Main results, mean $\pm$ std over 5 seeds. Exact match on the target task (adapt\_acc), per-token accuracy and exact match on the pretraining task after adaptation (pretask\_acc, higher means less forgetting). Bold/underline: best/second best as computed by \texttt{rh table}, including the unadapted row; they do not imply a significant difference (LoRA $r=4$, $r=8$ and scratch tie at 0.997 on descending sort).}\label{tab:main}\resizebox{\columnwidth}{!}{\begin{tabular}{llcccc}
\toprule
Method & Task & adapt\_acc $\uparrow$ & adapt\_tok\_acc $\uparrow$ & pretask\_acc $\uparrow$ & $n$ \\
\midrule
Head only & sort\_desc & 0.012 $\pm$ 0.007 & 0.615 $\pm$ 0.015 & 0.000 $\pm$ 0.000 & 5 \\
Last block & sort\_desc & 0.965 $\pm$ 0.007 & 0.995 $\pm$ 0.001 & 0.000 $\pm$ 0.000 & 5 \\
From scratch & sort\_desc & 0.997 $\pm$ 0.001 & 1.000 $\pm$ 0.000 & 0.000 $\pm$ 0.000 & 5 \\
Full fine-tuning & sort\_desc & 0.996 $\pm$ 0.003 & 1.000 $\pm$ 0.000 & \underline{0.000 $\pm$ 0.000} & 5 \\
\textbf{LoRA r=4} & sort\_desc & \underline{0.997 $\pm$ 0.001} & \underline{1.000 $\pm$ 0.000} & 0.000 $\pm$ 0.000 & 5 \\
LoRA r=8 & sort\_desc & \textbf{0.997 $\pm$ 0.001} & \textbf{1.000 $\pm$ 0.000} & 0.000 $\pm$ 0.000 & 5 \\
LoRA r=1 & sort\_desc & 0.978 $\pm$ 0.007 & 0.997 $\pm$ 0.001 & 0.000 $\pm$ 0.000 & 5 \\
LoRA r=2 & sort\_desc & 0.994 $\pm$ 0.002 & 0.999 $\pm$ 0.000 & 0.000 $\pm$ 0.000 & 5 \\
No adaptation & sort\_desc & 0.000 $\pm$ 0.000 & 0.093 $\pm$ 0.015 & \textbf{0.999 $\pm$ 0.001} & 5 \\
\midrule
Head only & reverse & 0.000 $\pm$ 0.000 & 0.248 $\pm$ 0.009 & \underline{0.004 $\pm$ 0.004} & 5 \\
Last block & reverse & 1.000 $\pm$ 0.000 & 1.000 $\pm$ 0.000 & 0.001 $\pm$ 0.000 & 5 \\
From scratch & reverse & \underline{1.000 $\pm$ 0.000} & \underline{1.000 $\pm$ 0.000} & 0.001 $\pm$ 0.000 & 5 \\
Full fine-tuning & reverse & \textbf{1.000 $\pm$ 0.000} & \textbf{1.000 $\pm$ 0.000} & 0.001 $\pm$ 0.000 & 5 \\
\textbf{LoRA r=4} & reverse & 0.999 $\pm$ 0.002 & 1.000 $\pm$ 0.000 & 0.001 $\pm$ 0.000 & 5 \\
LoRA r=8 & reverse & 0.941 $\pm$ 0.080 & 0.992 $\pm$ 0.011 & 0.001 $\pm$ 0.000 & 5 \\
LoRA r=1 & reverse & 0.799 $\pm$ 0.447 & 0.821 $\pm$ 0.400 & 0.000 $\pm$ 0.000 & 5 \\
LoRA r=2 & reverse & 0.800 $\pm$ 0.447 & 0.883 $\pm$ 0.261 & 0.000 $\pm$ 0.000 & 5 \\
No adaptation & reverse & 0.001 $\pm$ 0.000 & 0.193 $\pm$ 0.005 & \textbf{0.999 $\pm$ 0.001} & 5 \\
\bottomrule
\end{tabular}
}\end{table}

\begin{table*}[t]\centering\caption{Parameters and accuracy by system, 5 seeds. ``\% full'': share of the 70,016 parameters. ``min'': worst seed. ``fail'': seeds with reversal accuracy below 0.5.}\label{tab:rank}\begin{tabular}{lrrccccr}
\toprule
System & Params & \% full & sort\_desc & min & reverse & min & fail \\
\midrule
Head only & \rhval{main/head-only/reverse/trainable_params/mean} & \rhval{cmp/main/head-only/reverse/trainable_params/inv_ratio:pct1} & \rhval{main/head-only/sort_desc/adapt_acc/mean:3} $\pm$ \rhval{main/head-only/sort_desc/adapt_acc/std:3} & \rhval{main/head-only/sort_desc/adapt_acc/min:3} & \rhval{main/head-only/reverse/adapt_acc/mean:3} $\pm$ \rhval{main/head-only/reverse/adapt_acc/std:3} & \rhval{main/head-only/reverse/adapt_acc/min:3} & 5/5 \\
Last block & \rhval{main/last-block/reverse/trainable_params/mean} & \rhval{cmp/main/last-block/reverse/trainable_params/inv_ratio:pct1} & \rhval{main/last-block/sort_desc/adapt_acc/mean:3} $\pm$ \rhval{main/last-block/sort_desc/adapt_acc/std:3} & \rhval{main/last-block/sort_desc/adapt_acc/min:3} & \rhval{main/last-block/reverse/adapt_acc/mean:3} $\pm$ \rhval{main/last-block/reverse/adapt_acc/std:3} & \rhval{main/last-block/reverse/adapt_acc/min:3} & 0/5 \\
From scratch & \rhval{main/from-scratch/reverse/trainable_params/mean} & \rhval{cmp/main/from-scratch/reverse/trainable_params/inv_ratio:pct1} & \rhval{main/from-scratch/sort_desc/adapt_acc/mean:3} $\pm$ \rhval{main/from-scratch/sort_desc/adapt_acc/std:3} & \rhval{main/from-scratch/sort_desc/adapt_acc/min:3} & \rhval{main/from-scratch/reverse/adapt_acc/mean:3} $\pm$ \rhval{main/from-scratch/reverse/adapt_acc/std:3} & \rhval{main/from-scratch/reverse/adapt_acc/min:3} & 0/5 \\
Full fine-tuning & \rhval{main/full-fine-tuning/reverse/trainable_params/mean} & 100.0 & \rhval{main/full-fine-tuning/sort_desc/adapt_acc/mean:3} $\pm$ \rhval{main/full-fine-tuning/sort_desc/adapt_acc/std:3} & \rhval{main/full-fine-tuning/sort_desc/adapt_acc/min:3} & \rhval{main/full-fine-tuning/reverse/adapt_acc/mean:3} $\pm$ \rhval{main/full-fine-tuning/reverse/adapt_acc/std:3} & \rhval{main/full-fine-tuning/reverse/adapt_acc/min:3} & 0/5 \\
LoRA r=1 & \rhval{main/lora-r-1/reverse/trainable_params/mean} & \rhval{cmp/main/lora-r-1/reverse/trainable_params/inv_ratio:pct1} & \rhval{main/lora-r-1/sort_desc/adapt_acc/mean:3} $\pm$ \rhval{main/lora-r-1/sort_desc/adapt_acc/std:3} & \rhval{main/lora-r-1/sort_desc/adapt_acc/min:3} & \rhval{main/lora-r-1/reverse/adapt_acc/mean:3} $\pm$ \rhval{main/lora-r-1/reverse/adapt_acc/std:3} & \rhval{main/lora-r-1/reverse/adapt_acc/min:3} & 1/5 \\
LoRA r=2 & \rhval{main/lora-r-2/reverse/trainable_params/mean} & \rhval{cmp/main/lora-r-2/reverse/trainable_params/inv_ratio:pct1} & \rhval{main/lora-r-2/sort_desc/adapt_acc/mean:3} $\pm$ \rhval{main/lora-r-2/sort_desc/adapt_acc/std:3} & \rhval{main/lora-r-2/sort_desc/adapt_acc/min:3} & \rhval{main/lora-r-2/reverse/adapt_acc/mean:3} $\pm$ \rhval{main/lora-r-2/reverse/adapt_acc/std:3} & \rhval{main/lora-r-2/reverse/adapt_acc/min:3} & 1/5 \\
LoRA r=4 & \rhval{main/lora-r-4/reverse/trainable_params/mean} & \rhval{cmp/main/lora-r-4/reverse/trainable_params/inv_ratio:pct1} & \rhval{main/lora-r-4/sort_desc/adapt_acc/mean:3} $\pm$ \rhval{main/lora-r-4/sort_desc/adapt_acc/std:3} & \rhval{main/lora-r-4/sort_desc/adapt_acc/min:3} & \rhval{main/lora-r-4/reverse/adapt_acc/mean:3} $\pm$ \rhval{main/lora-r-4/reverse/adapt_acc/std:3} & \rhval{main/lora-r-4/reverse/adapt_acc/min:3} & 0/5 \\
LoRA r=8 & \rhval{main/lora-r-8/reverse/trainable_params/mean} & \rhval{cmp/main/lora-r-8/reverse/trainable_params/inv_ratio:pct1} & \rhval{main/lora-r-8/sort_desc/adapt_acc/mean:3} $\pm$ \rhval{main/lora-r-8/sort_desc/adapt_acc/std:3} & \rhval{main/lora-r-8/sort_desc/adapt_acc/min:3} & \rhval{main/lora-r-8/reverse/adapt_acc/mean:3} $\pm$ \rhval{main/lora-r-8/reverse/adapt_acc/std:3} & \rhval{main/lora-r-8/reverse/adapt_acc/min:3} & 0/5 \\
\bottomrule
\end{tabular}
\end{table*}

\begin{figure}[t]\centering\includegraphics[width=\columnwidth]{acc_vs_params.pdf}\caption{Exact-match accuracy of every seed against trainable parameters (log scale). Orange: LoRA, blue: other systems. Labels give the rank or system.}\label{fig:params}\end{figure}

\textbf{H1 (LoRA $r=4$ matches full fine-tuning with $<15\%$ of parameters): supported in this setting for rank 4 (the registered statement says rank $\ge4$; rank 8 on reversal (0.941, Table~\ref{tab:rank}) is below full fine-tuning (1.000) by more than the two-point margin, because of its seed failures, so H1 is not shown for rank 8 at the shared rate).} LoRA $r=4$ trains 7,168 parameters (\rhval{cmp/main/lora-r-4/sort_desc/trainable_params/inv_ratio:pct1}\% of 70,016). Its exact match is 0.997 against 0.996 for full fine-tuning on descending sort and 0.999 against 1.000 on reversal (Table~\ref{tab:main}); the paired differences (LoRA minus full fine-tuning) are \rhval{cmp/main/full-fine-tuning/sort_desc/adapt_acc/delta:4} ($p=$\rhval{cmp/main/full-fine-tuning/sort_desc/adapt_acc/paired_p:3}) and \rhval{cmp/main/full-fine-tuning/reverse/adapt_acc/delta:4} ($p=$\rhval{cmp/main/full-fine-tuning/reverse/adapt_acc/paired_p:3}, Table~\ref{tab:tests}). Both are far inside the two-point margin, but five seeds cannot rule out small differences, and both systems are at the ceiling, so the task discriminates poorly between them. A model trained from scratch with the same budget also reaches 0.997 and 1.000, so at $N=2000$ these tasks do not test whether pretraining is useful.

\textbf{H2 (accuracy rises with rank): supported on descending sort, not interpretable on reversal.} On descending sort the means are 0.978, 0.994, 0.997, 0.997 for $r=1,2,4,8$; rank 1 is below rank 4 by \rhval{cmp/main/lora-r-1/sort_desc/adapt_acc/delta:3} (paired $p=$\rhval{cmp/main/lora-r-1/sort_desc/adapt_acc/paired_p:4}, \texttt{rh compare} with rank 4 as reference). Ranks 4 and 8 are equal to the displayed precision. On reversal the means are 0.799, 0.800, 0.999, 0.941 (Table~\ref{tab:rank}). These are driven by whole-run failures, not by gradual loss: one seed at $r=1$ and one at $r=2$ end near chance with training loss far above the others (Table~\ref{tab:rank}, ``fail''), and at $r=8$ three seeds are at \rhval{main/lora-r-8/reverse/adapt_acc/min:2}--\rhval{run/2bfc86354c/adapt_acc:2}. The other seeds at $r=1$ reach \rhval{run/63db846dff/adapt_acc:3}--\rhval{main/lora-r-1/reverse/adapt_acc/max:3}. Reversal is therefore solved by rank 1 whenever optimisation succeeds, and the apparent rank trend is an optimisation-stability effect at the learning rate selected on one validation seed; we did not isolate its cause, and it is not evidence that rank 1 lacks capacity. The LoRA rate was tuned at $r=4$ only, so the failures at ranks 1, 2 and 8 may be artifacts of reusing it, and the rank trend on reversal should not be read as a property of rank. The rank-1 versus rank-4 difference on reversal is not significant ($p=$\rhval{cmp/main/lora-r-1/reverse/adapt_acc/paired_p:2}).

\textbf{H3 (LoRA forgets less than full fine-tuning): not supported, and uninformative because of a floor effect.} Every adapted system, LoRA at all ranks included, ends with pretask exact match of 0.000 on the descending-sort runs and at most 0.004 on average elsewhere (Table~\ref{tab:main}); forgetting is 0.9995 for all of them on descending sort (identical across systems empirically, so no test is defined, Table~\ref{tab:tests}). The residual 1/2000 of retention (0.0005) on reversal and in zero-shot rows is one test input whose reversal equals its ascending sort, not retained skill. Fig.~\ref{fig:curves} shows that retention is already zero at the first evaluation after 50 steps for full fine-tuning, LoRA $r=1$ and $r=4$, and last-block tuning. The task token did not protect the old skill, and neither the low rank nor freezing most weights helped; since even head-only tuning, which cannot adapt, sits at the floor and the pretrained model never learned to condition on the token, this test cannot discriminate between methods and we make no claim that LoRA does or does not forget less. The learning-rate sweep (Table~\ref{tab:lr}) shows the same at every learning rate tried, including settings where adaptation itself was incomplete. We cannot say whether a shorter schedule, replay, or a larger model would change this.

\textbf{H4 (last block worse than LoRA): supported on descending sort only.} Last-block tuning trains 34,496 parameters (against 7,168 for LoRA $r=4$) yet reaches 0.965 against 0.997 (difference \rhval{cmp/main/last-block/sort_desc/adapt_acc/delta:3} in favour of LoRA, paired $p=$\rhval{cmp/main/last-block/sort_desc/adapt_acc/paired_p:4}). On reversal the two are equal (0.999 vs.\ 1.000). Head-only tuning (896 parameters) fails on both tasks (0.012 and 0.000): with frozen features the output layer cannot realise a different permutation of the input.

\begin{table}[t]\centering\caption{Paired $t$-tests over the five seeds (main group): the statistics of \texttt{rh compare} with LoRA $r=4$ as the reference (A). Retention comparisons are identical across systems empirically (every system at the floor), so no test is defined.}\label{tab:tests}\resizebox{\columnwidth}{!}{\begin{tabular}{lllllr}
\toprule
A & B & metric & task & mean A$-$B & paired $p$ \\
\midrule
LoRA r=4 & Full fine-tuning & adapt\_acc & sort\_desc & +0.0003 & 0.816 \\
LoRA r=4 & Full fine-tuning & adapt\_acc & reverse & -0.0007 & 0.374 \\
LoRA r=4 & Last block & adapt\_acc & sort\_desc & +0.0314 & 0.000793 \\
LoRA r=4 & Last block & adapt\_acc & reverse & -0.0005 & 0.546 \\
LoRA r=1 & LoRA r=8 & adapt\_acc & sort\_desc & -0.0190 & 0.00369 \\
LoRA r=1 & LoRA r=8 & adapt\_acc & reverse & -0.1419 & 0.523 \\
LoRA r=4 & From scratch & adapt\_acc & sort\_desc & +0.0000 & 1 \\
LoRA r=4 & From scratch & adapt\_acc & reverse & -0.0007 & 0.374 \\
LoRA r=4 & Full fine-tuning & pretask\_acc & sort\_desc & +0.0000 & n/a \\
LoRA r=4 & Full fine-tuning & pretask\_acc & reverse & +0.0000 & n/a \\
\bottomrule
\end{tabular}
}\end{table}

\textbf{Where the method fails.} (i) LoRA at the selected learning rate has whole-run failures on reversal, at ranks 1, 2 and 8 but not 4 (Table~\ref{tab:rank}); $r=4$ was stable at that rate in all five seeds. (ii) With few examples, rank 8 collapses (Table~\ref{tab:ntrain}). (iii) Restricting LoRA to $\{q,v\}$ costs accuracy at low rank (Table~\ref{tab:targets}). (iv) No system retained the pretraining task, so the design cannot separate methods on forgetting. (v) Learning rates for full fine-tuning, LoRA, last block and head are at the upper grid edge.
